\subsection{自由李代数的构造}

定义 1. 集合 \(\mathbf{X}\) 上的自由李代数是商代数

\[
\mathrm{L} (\mathrm{X}) = \operatorname{Lib} (\mathrm{X}) / \mathrm{a},
\]

其中 a 是 \(\operatorname{Lib}(\mathbf{X})\) 的由具有下列形式之一的元素生成的双边理想

\begin{enumerate}
\def\labelenumi{(\arabic{enumi})}
\tightlist
\item
\end{enumerate}

\[
\mathrm{Q} (a) = a. a \text {   for   } a \text {   in   } \operatorname{Lib} (\mathrm{X}),\tag{2}
\]

\[
J (a, b, c) = a. (b. c) + b. (c. a) + c. (a. b)
\]

其中 \(a, b, c\) 属于 \(\operatorname{Lib}(\mathbf{X})\)。

显然 \(\mathrm{L}(\mathrm{X})\) 是李 K-代数；两个元素 \(u, v\)（\(\mathrm{L}(\mathrm{X})\) 的元素）的合成将记作 \([u, v]\)。当需要指明环 \(\mathbf{K}\) 时，我们用 \(\mathrm{L}_{\mathbb{K}}(\mathbf{X})\) 记 \(\mathrm{L}(\mathbf{X})\)。

下述命题说明了赋予 L(X) 自由李代数这一名称的理由。

命题 1. 设 \(\psi\) 为 \(\operatorname{Lib}(\mathbf{X})\) 到 \(\operatorname{L}(\mathbf{X})\) 上的典范映射，\(\phi\) 为 \(\psi\) 在 \(\mathbf{X}\) 上的限制。对每个映射 \(f\)（\(\mathbf{X}\) 到李代数 \(\mathfrak{g}\) 的映射），存在且仅存在一个同态 \(\mathrm{F}:\mathrm{L}(\mathrm{X})\to \mathfrak{g}\)，使得 \(f = \mathrm{F}\circ \phi\)。

\begin{enumerate}
\def\labelenumi{(\alph{enumi})}
\item
  F 的存在性：设 h 为 \(\operatorname{Lib}(\mathbf{X})\) 到 g 的扩张 \(f\) 的同态（第 1 号）。对 \(\operatorname{Lib}(\mathbf{X})\) 中的一切 a，\(h(\mathbf{Q}(a)) = h(a.a) = [h(a), h(a)] = 0\)； 类似地，g 所满足的 Jacobi 恒等式蕴含 \(h(\mathbf{J}(a, b, c)) = 0\)（对 \(\operatorname{Lib}(\mathbf{X})\) 中的 a, b, c）。 由此得出 \(h(a) = 0\)，从而存在 \(\operatorname{L}(\mathbf{X})\) 到 g 的同态 F，使得 \(h = F \circ \phi\)。 限制到 X 上，即得 \(f = F \circ \phi\)。
\item
  F 的唯一性：设 \(\mathrm{F}^{\prime}:\mathrm{L}(\mathrm{X})\to\mathfrak{g}\) 为使 \(f=\mathrm{F}^{\prime}\circ\phi\) 的同态。同态 \(F\circ\psi\) 与 \(F^{\prime}\circ\psi\)（\(\operatorname{Lib}(X)\) 到 g 的）在 X 上一致，因而相等；由于 \(\psi\) 是满射，\(F=F^{\prime}\)。
\end{enumerate}

推论 1. 族 \((\phi (x))_{x\in \mathbf{X}}\) 在 \(\mathbf{K}\) 上是 \(\mathrm{L}(\mathbf{X})\) 中的自由族。

设 \(x_{1}, x_{2}, \ldots, x_{n}\) 为 X 中互异的元素，\(\lambda_{1}, \ldots, \lambda_{n}\) 为 K 中满足下式的元素

\[
\lambda_ {1}. \phi (x _ {1}) + \dots + \lambda_ {n}. \phi (x _ {n}) = 0.\tag{3}
\]

设 \(\mathfrak{g}\) 为以 \(\mathbf{K}\) 为基础模的交换李代数。对 \(i = 1,2,\ldots ,n\)，存在同态 \(\mathrm{F}_i\)（\(\mathbf{L}(\mathbf{X})\) 到 \(\mathfrak{g}\) 的），使得 \(\mathrm{F_i}(\phi (x_i)) = 1\) 且 \(\mathrm{F_i}(\phi (x)) = 0\) 对 \(x\neq x_{i}\) 成立（命题 1）；把 \(\mathrm{F}_i\) 应用于关系 (3)，便得 \(\lambda_{i} = 0\)。

推论 2. 设 \(\mathfrak{a}\) 为李代数。\(\mathrm{L}(\mathrm{X})\) 被 \(\mathfrak{a}\) 的每个扩张都是非本质的。

设 \(\mathfrak{a} \xrightarrow{\lambda} \mathfrak{g} \xrightarrow{\mu} \mathrm{L}(\mathrm{X})\) 为这样一个扩张（第 I 章，第 1 节，第 7 号）。由于 \(\mu\) 是满射，存在映射 \(f\)（\(\mathrm{X}\) 到 \(\mathfrak{g}\) 的），使得 \(\phi = \mu \circ f\)。 设 \(\mathrm{F}\) 为 \(\mathrm{L}(\mathrm{X})\) 到 \(\mathfrak{g}\) 的使 \(f = \mathrm{F} \circ \phi\) 的同态（命题 1）。 于是 \((\mu \circ \mathrm{F}) \circ \phi = \mu \circ f = \phi\)，命题 1 表明 \(\mu \circ \mathrm{F}\) 是 \(\mathrm{L}(\mathrm{X})\) 的恒同自同构。 因此所给的扩张是非本质的（第 I 章，第 1 节，第 7 号，命题 6 与定义 6）。

由于环 K 非零，命题 1 的推论 1 表明 \(\phi\) 是单射。因此集合 X 可借助 \(\phi\) 等同于它在 L(X) 中的像；在这一约定下，X 生成 L(X)，且 X 到李代数 \(\mathfrak{g}\) 的每个映射都可扩张为 L(X) 到 \(\mathfrak{g}\) 的李代数同态。

注. 当 X 为空时，M(X) 为空，因而 L(X) = \{0\}。若 X 由单个元素 x 组成，则 L(X) 的子模 K.x 是

\(L(X)\) 的李子代数；由于 X 生成 \(L(X)\)，命题 1 的推论 1 表明 \(L(X)\) 是以 \(\{x\}\) 为基的自由模。

